\documentclass[11pt,a4paper]{article}
\usepackage[margin=1in]{geometry}
\usepackage{amsmath,amssymb,amsthm}
\usepackage{booktabs}
\usepackage{xcolor}
\definecolor{linkcol}{rgb}{0.05,0.3,0.45}
\usepackage[colorlinks=true,linkcolor=linkcol,citecolor=linkcol,urlcolor=linkcol]{hyperref}

\newtheorem{theorem}{Theorem}
\newtheorem{proposition}{Proposition}
\newtheorem{lemma}{Lemma}
\newtheorem{corollary}{Corollary}
\theoremstyle{definition}
\newtheorem{example}{Example}
\newtheorem{remark}{Remark}
\newtheorem{conjecture}{Conjecture}

\newcommand{\ket}[1]{\lvert #1 \rangle}
\newcommand{\bra}[1]{\langle #1 \rvert}
\newcommand{\proj}[1]{\lvert #1 \rangle\!\langle #1 \rvert}
\newcommand{\tr}{\operatorname{Tr}}
\newcommand{\id}{\mathbb{1}}

\title{The Nonclassicality Network: a signed theorem graph for quantum
foundations with machine-visible open problems}

% TODO (author): fill in name(s), affiliation and email before submission.
% Consider disclosing AI-assisted drafting/curation if required by the venue's policy;
% the curation and verification protocol is described in Sec. 6.
\author{[Author name]\\ \small [Affiliation] \quad [email]}
\date{October 2026 \quad (v2.4 of the network)}

\begin{document}
\maketitle

\begin{abstract}
We present a curated, typed, signed theorem graph of quantum foundations,
quantum information, quantum thermodynamics and resource theories, released as
an open repository with two interactive pages and machine-readable data. Nodes
are typed (axioms, models, state-operator clusters, process-operator clusters,
witnesses, effects, tasks); edges are signed relations carrying verbatim
statements and references, with impossibility (no-go) theorems recorded as
first-class negative edges, so that the \emph{absence} of an edge always means
``unsettled'' rather than ``false''. Witnesses are operator-valued nodes paired
one-to-one with the effects they detect, each carrying structured requirement
and blindness records; a state$\times$process verdict matrix then enumerates
open cells mechanically (in the current release, $445$ of $3094$ cells are
settled). We demonstrate that the structure is generative: closing formerly
uncharted cells yields several propositions proved here, including a
dephasing-factorization theorem (covariant processes with time-translation-symmetric
effects form simplex-embeddable, hence noncontextual, fragments for every input
state), a classification of time-translation-symmetric bipartite correlations
(for nondegenerate noninteracting Hamiltonians all such states are
classical--classical; quantum correlations survive only in degeneracy blocks),
a composed corollary exhibiting Bell-local yet contextual scenarios from
Werner states, and independence counterexamples separating magic from discord.
Fifteen named open problems are registered. The intended use is threefold:
semi-automated generation of research problems from the gap census; a
reviewable survey format in which checking the review means auditing
individually cited edges; and an interactive visualization of the field's
logical structure. All sources, data and pages are provided for community
extension.
\end{abstract}

\section{Introduction}

Quantum foundations is unusually rich in \emph{impossibility} results: Bell's
theorem, the Kochen--Specker theorem, no-cloning and no-broadcasting, the PBR
theorem, and the modern no-go theorems of generalized contextuality. It is
equally rich in \emph{possibility} and \emph{equivalence} results: the
hierarchy of Bell nonlocality, steering and entanglement
\cite{wiseman2007}, the equivalences between nonclassicality, incompatibility,
no-broadcasting, superposition and entangleability in generalized
probabilistic theories (GPTs)
\cite{plavala2016,barnum2007,aubrun2021,aubrun2022}, and the equivalences
tying Wigner negativity to contextuality in odd prime dimensions
\cite{howard2014} and in continuous variables \cite{booth2022}. These results
are scattered across subfields that use different formalisms, and simple
questions of the form ``does property $X$ imply property $Y$, in which
fragment, for which states and which measurements?'' are surprisingly hard
to answer from the literature.

This paper introduces a working tool addressed at exactly that difficulty: a
\emph{signed theorem graph} --- the Nonclassicality Network --- whose design
goal is that missing knowledge be visible. Three design decisions do the work.

\begin{enumerate}
  \item \textbf{Three-valued edges.} Between any two concepts the graph
  distinguishes (i) a positive edge (a theorem asserting the relation), (ii) a
  negative edge (a theorem asserting the relation \emph{cannot} hold: a no-go,
  an exclusion, a counterexample), and (iii) no edge, which by convention means
  \emph{unsettled}. Without (ii), impossibility theorems would be encoded as
  absences and become indistinguishable from ignorance; a search for open
  problems would then re-propose closed questions. Negative edges are drawn
  with inhibition-style flat-bar terminators ($\dashv$) and are filterable.

  \item \textbf{Typed layers with operator-valued witnesses.} Nodes are typed:
  axioms (including a ``genesis'' family: symplectic structure, epistemic
  restriction, the Poisson$\to$Moyal lift, the complex unit $i$, and $\hbar$),
  models, state-operator clusters, process-operator clusters, witnesses,
  effects, and tasks. Every operational effect is paired one-to-one with the
  witness (functional, inequality, semidefinite or linear program, protocol)
  that detects it --- the design principle \emph{no effect without a witness}
  --- and each witness node carries structured records: what it requires at
  the model, state and process level, and on which clusters it is provably
  blind. The recorded exception to strict pairing is threshold stratification:
  the CHSH functional detects Bell violation, the Tsirelson boundary and PR
  correlations at thresholds $2$, $2\sqrt2$, $4$
  \cite{chsh1969,tsirelson1980,pr1994}.

  \item \textbf{Mechanical gap enumeration.} For each witness, a
  state-cluster$\times$process-cluster matrix is populated from recorded
  theorems; every unfilled cell is, by construction, an open question. The
  current release settles $445$ of $3094$ cells; the remaining $2649$ cells,
  together with $15$ named open problems, constitute a machine-generated
  research agenda.
\end{enumerate}

\paragraph{Motivation and intended use.}
The network was built for three purposes, and the design decisions above map
onto them directly.
\emph{(i)~Automating research-idea generation.} Because the absence of an edge
means ``unsettled'' and never ``false'', the set of open questions is a
queryable object rather than an impression of the literature: every unfilled
matrix cell and every missing edge between typed nodes is a candidate theorem,
pre-sorted by the structure into four closure routes --- compose recorded
edges, prove directly, exhibit a counterexample, or locate the result in the
literature. Section~\ref{sec:generated} demonstrates that this loop produces
nontrivial output: thirteen propositions across three curation rounds, each of
which began as a specific uncharted cell.
\emph{(ii)~Reviewing a field as curation.} The graph is a review article in a
different format: each claim of the review is a single edge carrying a
verbatim statement, references and a provenance status, so checking the review
means auditing edges one at a time, and disputing it means objecting to a
specific edge identifier rather than to a paragraph. Releases are versioned
and each records its gap census, so the review does not silently decay.
\emph{(iii)~Visualizing the field.} The layered, signed, filterable pages make
the logical shape of nonclassicality visible at a glance --- where possibility
flows, where no-go bars block, which clusters a given witness can never see
--- and the witness panels answer the operational question ``what is required,
at model, state and process level, to observe effect $X$?'' interactively.
The review process is itself the generator: surveying the field in this format
yields, as a by-product, its own research agenda.

The repository contains two interactive pages over the same data --- a flat
version (v1) and the layered version (v2) --- documenting the evolution of the
structure, together with machine-readable JSON carrying a verbatim statement
and references for every edge. Section~\ref{sec:backbone} summarizes the
curated backbone. Section~\ref{sec:generated} contains the generative payoff:
propositions obtained by closing formerly uncharted cells, with proofs.
Section~\ref{sec:open} lists the registered open problems.
Section~\ref{sec:protocol} states the curation and verification protocol ---
what was verified against the literature, what is cited from standard
knowledge, and what is new here --- and Section~\ref{sec:repo} describes the
repository.

\section{The curated backbone}\label{sec:backbone}

Two nested classicality boundaries organize the graph. First, for a single GPT
system the following are equivalent characterizations of
\emph{nonclassicality} (non-simpliciality of the state cone): the existence of
incompatible measurements \cite{plavala2016}; the impossibility of universal
broadcasting \cite{barnum2007}; the existence of ``superposition'' (strong
incompatibility) \cite{aubrun2022}; and \emph{entangleability}, i.e.\ the
strict separation of the minimal and maximal tensor products with any other
nonclassical system --- the resolution of Barker's conjecture
\cite{aubrun2021,aubrun2022}. Second, and strictly stronger, a prepare--measure
fragment admits a generalized-noncontextual ontological model
\cite{spekkens2005,spekkens2008} if and only if it is
\emph{simplex-embeddable} \cite{schmid2021}; the complete test is a linear
program \cite{selby2022lp}. The epistricted theories --- Spekkens' toy theory
\cite{spekkens2007,spekkens2016}, Gaussian quantum mechanics as epistricted
Liouville mechanics \cite{brs2012}, and the odd-dimensional stabilizer
subtheory with Gross's nonnegative Wigner function
\cite{gross2006,catani2017,schmid2022} --- live between the two boundaries:
they exhibit incompatibility, no-cloning, interference phenomenology and
steering-like ensembles while remaining noncontextual. The qubit stabilizer
subtheory is the known anomaly: it admits no nonnegative diagram-preserving
representation and is state-independently contextual
\cite{schmid2022,mermin1990,peres1991,cabello2008}.

Beyond this backbone the graph records (with per-edge references in the data
files): the nonlocality--steering--entanglement hierarchy and its strictness
\cite{wiseman2007,werner1989,barrett2002,gisin1991}; incompatibility as
necessary for Bell violations \cite{fine1982} and equivalent to steerability
\cite{quintino2014,uola2014}; the failure of the Peres conjecture
\cite{vertesi2014}; the negativity--contextuality equivalences
\cite{howard2014,booth2022} with the representation-relativity caveat of the
nonnegative $\Lambda$-polytope model \cite{zurel2020}; contextuality as a
precondition for entanglement \cite{plavala2024}; the conditional equivalence
of incompatibility and contextuality when all states are available
\cite{tavakoli2020} against its unconditional failure \cite{selby2023};
no-broadcasting as a characterization of operational contextuality
\cite{nobroadcast2024}; thermodynamic witnesses of contextuality
\cite{lostaglio2018,perarnau2017,micadei2019}; coherence as asymmetry and its
thermodynamic role \cite{marvian2014,lostaglio2015}; the genesis decomposition
(Groenewold--van Hove obstruction and the Moyal deformation
\cite{groenewold1946,moyal1949,bffls1978,kontsevich2003}; Hudson's theorem
\cite{hudson1974}; non-quadratic generators as the source of continuous-variable
universality and negativity \cite{lloyd1999,takagi2018}; local tomography and
the complex unit \cite{hardywootters2012,renou2021,stueckelberg1960}); the
almost-quantum obstruction to the no-restriction hypothesis
\cite{navascues2015,sainz2018}; and the Sorkin hierarchy
\cite{sorkin1994,sinha2010,dakic2014}.

\section{Propositions generated by closing uncharted cells}\label{sec:generated}

Each result below occupied a cell or node pair with no recorded connecting or
disconnecting theorem. We state the network provenance (which edges were
composed) with each item. Results labelled ``derived'' are proved here; we do
not claim they are absent from the literature, only that we did not find them
stated (see Sec.~\ref{sec:protocol}).

\subsection{Round 1: boundary corollaries}

\begin{proposition}[Contextuality implies entangleability; provenance
$A5\circ A3$]\label{prop:p1}
If a tomographically complete prepare--measure fragment of a GPT system $A$
fails generalized noncontextuality, then for every nonclassical system $B$
there exists a consistent composite containing entangled states or entangled
effects.
\end{proposition}
\begin{proof}
Contextuality implies failure of simplex embeddability \cite{schmid2021},
hence the state cone of $A$ is not simplicial. By the entangleability of cones
\cite{aubrun2021,aubrun2022}, the minimal and maximal tensor products of $A$
with any nonclassical $B$ differ, and any composite beyond the minimal one
contains entangled elements.
\end{proof}

The converse fails (the toy theory is nonsimplicial yet noncontextual
\cite{spekkens2007,schmid2021}), and the unique escape is postulating the
minimal tensor product: \emph{separable quantum theory} retains Wigner
negativity, magic and contextuality of its quantum marginals while possessing
no entangled, steerable or Bell-violating states, and purification
\cite{cdp2010} fails in it by construction (cones between the minimal and
maximal tensor product are standard \cite{janotta2014}). A complementary
necessity direction is the theorem of Pl\'avala and G\"uhne \cite{plavala2024}.
Similarly, a fragment whose states and effects are diagonal in one basis is
simplex-embeddable, so basis-independent coherence is necessary for
contextuality, while epistricted counterexamples
\cite{spekkens2007,catani2023} and vanishing-coherence contextuality
\cite{rossi2023} show it is not sufficient and has no threshold; quantitative
scenario-specific witnesses exist \cite{wagner2024}.

\subsection{Round 2: time-translation symmetry hides contextuality}

Let $H$ be a Hamiltonian with nondegenerate spectrum, $D$ the associated
energy-dephasing map, and call a state, effect or channel \emph{TTS}
(time-translation symmetric, covariant) when it commutes with the unitary
group $U_t=e^{-iHt}$.

\begin{theorem}[Covariant fragments are noncontextual; provenance: uncharted
cell TTS$\times$covariant]\label{thm:n1}
Consider any prepare--transform--measure fragment whose transformations $T$
are covariant and whose effects $E$ are TTS. Then for nondegenerate $H$ the
fragment's statistics satisfy
$\tr\!\big(E\,T(\rho)\big)=\tr\!\big(E\,T(D(\rho))\big)$ for every state
$\rho$; the fragment factors through the commutative energy-diagonal
subtheory, is simplex-embeddable and hence noncontextual, its Kirkwood--Dirac
distributions over energy chains are genuine probabilities, and no coherence
monotone is estimable from it.
\end{theorem}
\begin{proof}
For TTS effects, $U_t^\dagger E U_t=E$, so
$\tr(E\,T(U_t\rho U_t^\dagger))=\tr(U_t^\dagger E U_t\,T(\rho))=\tr(E\,T(\rho))$
using covariance of $T$. Averaging over $t$ yields invariance under
pre-dephasing. For nondegenerate $H$, $D$ maps onto the diagonal (commutative)
algebra, TTS effects are diagonal, and covariant channels preserve the
diagonal subalgebra, so the fragment is tomographically equivalent to a
classical stochastic fragment, which embeds in a simplex with its dual
effects; noncontextuality follows from the characterization of
\cite{schmid2021}.
\end{proof}

Theorem~\ref{thm:n1} says TTS processes do not lack contextuality but
\emph{hide} it: the hiding map is exactly $D$. Lostaglio's noncontextuality of
two-point-measurement fluctuation theorems \cite{lostaglio2018} is the special
case $T=\mathrm{id}$, $E$ the energy projectors, and the no-go of
\cite{perarnau2017} is the statement that the first projective energy
measurement implements the dephasing. For degenerate $H$ the factorization is
only block-wise and covariant fragments can remain contextual (open problem
SG1). Conversely, the state side carries no blindness at all:

\begin{proposition}[No state-blind cluster for contextuality]\label{prop:n2}
Every quantum state --- Gibbs states and the maximally mixed state included ---
violates state-independent contextuality inequalities (the Peres--Mermin
inequality of \cite{cabello2008}; the Yu--Oh 13-ray inequality
\cite{yuoh2012}). Consequently, for nondegenerate $H$, any state-independent
contextuality witness set necessarily contains non-TTS measurements.
\end{proposition}
\begin{proof}
The first statement is the content of state-independence in
\cite{cabello2008,yuoh2012}. The second follows from Theorem~\ref{thm:n1}:
a TTS witness set would yield a noncontextual fragment for every state.
\end{proof}

\subsection{Round 3: new closures}

\begin{proposition}[Symmetric bipartite states of generic noninteracting pairs
are classically correlated; provenance: uncharted cell
TTS$\times$correlation clusters]\label{prop:r1}
Let $H=H_A+H_B$ with nondegenerate total spectrum (all sums $E^A_i+E^B_j$
distinct). Then every TTS state is classical--classical: zero discord, hence
separable, unsteerable and Bell-local. Entangled or discordant TTS states
exist iff the total spectrum is degenerate, and all their quantum correlations
are supported inside degeneracy blocks.
\end{proposition}
\begin{proof}
Nondegeneracy of the total spectrum forces the eigenbasis of $H$ to be the
unique product basis $\ket{E^A_i}\otimes\ket{E^B_j}$ (any local degeneracy or
cross-degeneracy would produce a degenerate sum). A state commuting with a
nondegenerate $H$ is diagonal in its eigenbasis, hence of the form
$\sum_{ij}p_{ij}\proj{ij}$, a classical--classical state: zero discord
\cite{ollivier2001,henderson2001}, separable, and therefore unsteerable and
Bell-local \cite{wiseman2007}. For the degenerate case, two equal-gap qubits
admit the TTS maximally entangled state
$(\ket{01}+\ket{10})/\sqrt2$ supported in the degenerate $\{01,10\}$ block.
\end{proof}

\begin{remark}
Proposition~\ref{prop:r1} is elementary and may be folklore; we did not find
it stated as a classification. It is consistent with experiment: the
correlation-driven reversal of heat flow of Micadei \emph{et al.}
\cite{micadei2019} employs initial coherence $\ket{01}\!\bra{10}$ --- exactly
a degenerate-block term for equal-frequency qubits, the sole habitat the
proposition allows.
\end{remark}

\begin{lemma}[TTS measurements are compatible]\label{lem:r2}
For nondegenerate $H$, any two TTS sharp observables commute (each is a
function of $H$), and any TTS POVM is a classical post-processing of the
energy measurement; the whole TTS measurement cluster is jointly measurable.
Consequently TTS probes admit no incompatibility, hence no steering
\cite{quintino2014,uola2014} and no Bell violation \cite{fine1982}, and the
commutativity of the fragment reproves the blindness column of
Theorem~\ref{thm:n1}.
\end{lemma}
\begin{proof}
$[A,H]=0$ with $H$ nondegenerate forces $A$ diagonal in the energy basis; two
diagonal operators commute. A POVM with diagonal effects
$E_k=\sum_i e_{k|i}\Pi_i$ is simulated by measuring $H$ and sampling $k$ from
$e_{k|i}$.
\end{proof}

\begin{corollary}[Bell-local yet contextual scenarios; provenance
$A8\circ B3$]\label{cor:r3}
There exist bipartite entangled states admitting local hidden variable models
for all single-copy measurements \cite{werner1989,barrett2002} whose derived
sequential (remote-state-preparation) scenario, equipped with all measurements
on the steered side, is preparation-and-measurement contextual. Generalized
contextuality is therefore a strictly finer detector of composite
nonclassicality than Bell violation.
\end{corollary}
\begin{proof}
Theorem~1 of \cite{plavala2024} states that a preparation-and-measurement
noncontextual model for the steered states together with all measurements
implies separability. The contrapositive applies to every entangled state,
in particular to Werner states inside the LHV region. Bell locality for all
single-copy measurements is Refs.~\cite{werner1989,barrett2002}.
\end{proof}

\begin{remark}
The corollary is an immediate inference from \cite{plavala2024}; we verified
that it is not stated there explicitly. Note the fragment scope: the
contextual fragment contains \emph{all} measurements on the steered side;
restricted fragments can remain noncontextual --- e.g.\ entangled stabilizer
states probed by stabilizer measurements admit the local nonnegative Gross
model \cite{gross2006,schmid2022}. There is no tension between the two
statements.
\end{remark}

\begin{example}[Magic and discord are logically independent]\label{ex:r4}
(a)~$\ket{T}\otimes\ket{0}$, with $\ket{T}$ a canonical non-stabilizer state
\cite{bravyi2005}, is a pure product state --- zero discord, zero entanglement
--- yet magic. (b)~$\rho=\tfrac12(\proj{0}\otimes\proj{0}+\proj{+}\otimes\proj{+})$
is a mixture of stabilizer product states (hence inside the stabilizer
polytope, non-magic) but has nonzero discord on both sides: if
$[\rho,A\otimes\id]=0$ then $[P_0,A]\otimes P_0=-[P_+,A]\otimes P_+$, and since
$P_0,P_+$ are linearly independent both commutators vanish; $A$ then commutes
with $P_0$ and $P_+$, which generate the full qubit algebra, so
$A\propto\id$. Thus no nontrivial local observable commutes with $\rho$, which
therefore is not classical on either side \cite{ollivier2001,henderson2001}.
Together (a) and (b) close the formerly uncharted magic--discord pair in both
directions: no unconditional edge can connect them.
\end{example}

\begin{proposition}[Sorkin order of epistricted theories;
inheritance]\label{prop:r6}
Gaussian quantum mechanics and the odd-dimensional stabilizer subtheory
satisfy the Sorkin truncation $I_3=0$ identically.
\end{proposition}
\begin{proof}
Both are subtheories of quantum theory --- their states, transformations and
measurements are quantum objects with quantum statistics
\cite{brs2012,gross2006} --- and $I_3$ vanishes identically in quantum theory
\cite{sorkin1994}.
\end{proof}

\begin{conjecture}[SG4]
The discrete toy (field) theory, which is \emph{not} a quantum subtheory
(it does not reproduce GHZ-type stabilizer statistics
\cite{spekkens2007,catani2017}), has an undetermined Sorkin order: define its
three-path interferometer and compute $I_3$. Either outcome is informative.
\end{conjecture}

Finally, one formerly uncharted cell is closed by existing literature that the
graph had not yet absorbed: the Kirkwood--Dirac analogue of Hudson's boundary.
Pure KD-classical states are characterized by support-uncertainty
(``complete incompatibility'') conditions \cite{debievre2021}; noncommutation
alone does not suffice for KD nonpositivity \cite{ady2021}; and for discrete-Fourier
bases a pure state is KD-classical iff $n_A(\psi)\,n_B(\psi)=d$ \cite{xu2022}.

\section{Registered open problems}\label{sec:open}

The repository registers fifteen named problems; abbreviated here. G1:~a
steering analogue of the entangleability of cones (with a stated conjecture).
G2:~quantitative coherence$\to$contextuality bounds, approachable by sweeping
fragments with the simplex-embedding linear program \cite{selby2022lp}.
G3:~an even-dimensional magic--contextuality equivalence, obstructed by
state-independent Pauli contextuality and by the $\Lambda$-polytope model
\cite{zurel2020}. G4:~completeness of thermodynamic contextuality witnesses
within thermal operations (sharpened by Theorem~\ref{thm:n1}: probes must
leave the TTS cluster). G5:~contextuality quantifiers for higher-order
interference theories \cite{dakic2014}. G6:~which thermodynamic requirements
force purification-type axioms. G7:~a monogamy calculus between contextuality
and nonlocality. G8:~NPT bound entanglement and its thermodynamics.
G9:~a principled measurement-restriction rule for almost-quantum theories
\cite{navascues2015,sainz2018}. G10:~a complexity-theoretic column for every
witness (contextuality is a linear program; separability is NP-hard
\cite{gurvits2003}; quantum-set membership is undecidable \cite{jnvwy2020}).
SG1:~covariant contextuality inside degeneracy blocks (the failure mode of
Theorem~\ref{thm:n1}). SG2:~threshold stratification of witness families.
SG3:~the uncharted-cell program itself ($2649$ cells at this release).
SG4:~the toy-field-theory Sorkin order (above). SG5:~hidden-contextuality
activation: do local filters, which activate hidden nonlocality
\cite{werner1989,barrett2002}, amplify the sequential-scenario contextuality
of Corollary~\ref{cor:r3}?

\section{Curation and verification protocol}\label{sec:protocol}

Reliability is handled by explicit provenance rather than by implication.
(i)~Every edge in the data files carries a verbatim statement and references;
the per-edge reference list, not this paper, is the citation of record for
compiled results. (ii)~Edges marked \emph{theorem} are literature claims; the
load-bearing post-2019 results were verified against bibliographic sources
during curation (October 2026), while classic results (e.g.\
\cite{hudson1974,groenewold1946,moyal1949,fine1982,werner1989}) are cited from
standard knowledge and their page data should be rechecked before quotation.
(iii)~Edges marked \emph{derived} are proved in this paper or in the
repository catalogs and are \emph{not} claimed to be previously unpublished;
Propositions~\ref{prop:p1}, \ref{prop:r1} and Lemma~\ref{lem:r2} are
elementary enough to be folklore, and a dedicated novelty pass is advised
before attributing priority. (iv)~Composed corollaries state exactly which
recorded theorems they compose, so a reader can audit each in one step.
(v)~The authors take responsibility for remaining errors and welcome
corrections as repository issues; the graph is versioned, and every release
records its gap census.

\section{Repository}\label{sec:repo}

The repository (\url{https://github.com/tahmeedhossain/GEN_GPT_Level}, served as live pages at \url{https://tahmeedhossain.github.io/GEN_GPT_Level/}) contains:
\texttt{index.html} (overview and the generated-proposition record);
\texttt{v1.html} (flat network, 81 edges); \texttt{v2.html} (layered network,
148 signed edges, witness records, verdict matrix with uncharted-cell census);
\texttt{data/network\_v1.json} and \texttt{data/network\_v2.json}
(machine-readable: nodes, signed edges with statements and references, witness
records, combinations, polarity metadata); \texttt{catalog/} (prose catalogs
with full proofs); and this paper's source under \texttt{paper/}. The data
schema is documented in the README; contributions of new edges are expected to
carry a statement, a reference or proof, a kind, a status and a polarity.

\paragraph{Acknowledgements.} % TODO (author): add acknowledgements/funding as appropriate.

\bibliographystyle{unsrt}
\bibliography{refs}

\end{document}
